% Converted from the authoritative recovery source, 2026-11-14.
\subsection{Posterior-access transfer and generalized influence on an
actual
law}\label{influence-recovery:posterior-access-transfer-and-generalized-influence-on-an-actual-law}

Suppose a coupling \((X,Z)\) has \(X\sim\mu\) and a product-block
marginal \(\nu\) for \(Z\), with blocks of size at most \(k\). Let \(F\)
map the entire source cube to the physical cube. Assume:

\begin{enumerate}
\def\labelenumi{\arabic{enumi}.}
\tightlist
\item
  Each output \(F_i\) has a deterministic decision tree making at most
  \(r\) source-block requests. Its block degree is consequently at most
  \(r\).
\item
  Conditional on every supported \(x\),
  \(\Pr[F(Z)_i\ne x_i\mid X=x]\le\gamma\) for every \(i\).
\item
  An arbitrary adaptive sequence of \(s\) source-block requests can be
  supplied consistently from \(Z\mid X=x\) using at most \(B(s)\)
  physical-coordinate queries.
\end{enumerate}

Let \(p\) be globally \([0,1]\)-valued of degree \(d\). Set
\(g=p\circ F\) and \(\Delta=rd\). Its block degree is at most
\(\Delta\). Define the recoverable-source influence

\[
 \mathfrak J_a(p;F,\nu)
 :=\frac12\mathbb E_{Z,Z'_a}
       [p(F(Z))-p(F(Z^{a\leftarrow Z'_a}))]^2.
 \tag{S1}
\]

If \(v=\operatorname{Var}_\mu(p)>0\) and \(Cd^2\gamma\le v/4\), (RV) and
(B2) prove

\[
 \boxed{\quad
 \max_a\mathfrak J_a(p;F,\nu)
 \ge c'\frac{v^{5/2}}
 {(rd)^6[k+\log(4rd/v)]^6}.
 \quad} \tag{S2}
\]

These are collective counterfactual changes generated by resampling one
independent cause. For the extreme ordered law \(X=(Z,\ldots,Z)\) with
one fair spin \(Z\), the magnetization has ordinary normalized flip
influence \(n^{-2}\), while its source-block resampling influence is
one; observing any physical spin reveals the source. (For the
\([0,1]\)-valued acceptance \((1+M)/2\) these numbers are respectively
\(1/(4n^2)\) and \(1/4\).) A cause can affect many physical sites, which
is why (S2) evades the fixed-size physical-update obstruction. For an
approximate \(F\), (S1) compares two outputs of \(F\); it must not be
described as the physical single-spin influence under \(\mu\).

For \(e\in(0,1)\), choose \(\gamma\le e/(2Cd^2)\). Apply the source tree
from the preceding potential argument to \(g\) with error \(e/2\) and
answer its source-block queries using the posterior oracle in item 3.
The tree ignores all additionally observed physical data.
Unconditionally, its supplied source answers have exactly the product
law \(\nu\), so its original potential analysis and \(e/2\)
approximation guarantee remain valid. By (R), its error relative to
\(p(X)\) is at most \(e\). The number of physical queries is at most

\[
 B\bigl(S(rd,e/2;k)\bigr).
 \tag{S3}
\]

Implement the posterior kernels using preassigned independent uniform
random seeds. Averaging fixes those seeds to give a deterministic
physical decision tree with the same average error and no greater depth.
Extra physical observations must not be fed back into the source-block
selection rule: its analysis is under the source transcript, rather than
a claim that source independence survives arbitrary physical
conditioning.
